\documentclass{umk-matematika}
\begin{document}

\begin{center}
{\Large\bfseries Практическая работа 24}
\end{center}
\vspace{0.5cm}

\textbf{Задача 1.} Решите уравнение: $\log\_{2} x = 3$\par\vspace{10pt}
\textbf{Задача 2.} Решите уравнение: $\log\_{5} (x - 1) = 2$\par\vspace{10pt}
\textbf{Задача 3.} Решите уравнение: $\log\_{3} (3x) = 2$\par\vspace{10pt}
\textbf{Задача 4.} Решите уравнение: $\log\_{4} (x + 3) = 1$\par\vspace{10pt}
\textbf{Задача 5.} Решите уравнение: $\log\_{2} (x + 1) + \log\_{2} (x - 1) = 3$\par\vspace{10pt}
\textbf{Задача 6.} Решите уравнение: $\log\_{3} (x^2 - 4) - \log\_{3} (x - 2) = 1$\par\vspace{10pt}
\textbf{Задача 7.} Решите уравнение: $\log\_{2}^{2} x - 3\log\_{2} x + 2 = 0$\par\vspace{10pt}
\textbf{Задача 8.} Срок службы оборудования $t$ (лет) связан с износом $I$ (\%) формулой $t = 5 \cdot \log\_{2} \left(\frac{100}{100 - I}\right)$. Через сколько лет износ достигнет $75\\%$?\par\vspace{10pt}
\textbf{Задача 9.} Решите уравнение: $\log\_{2} (x + 2) + \log\_{2} (x - 2) = \log\_{2} (5x - 10)$\par\vspace{10pt}
\textbf{Задача 10.} Решите уравнение: $\log\_{5} x + \log\_{x} 5 = 2$\par\vspace{10pt}
\newpage
\section*{Ответы}

\begin{enumerate}
  \item Ответ: $x = 8$
  \item Ответ: $x = 26$
  \item Ответ: $x = 3$
  \item Ответ: $x = 1$
  \item Ответ: $x = 3$
  \item Ответ: нет решений
  \item Ответ: $x = 2;\ x = 4$
  \item Ответ: через $10$ лет
  \item Ответ: $x = 3$
  \item Ответ: $x = 5$
\end{enumerate}
\end{document}
